\section{从二分类到 Softmax 分类}
\label{sec:13-Machine-Learning/02-Linear-Models/03}

你刚学会线性回归，兴致勃勃地拿它去预测客户会不会续费。响应变量 \(y\) 取 0 和 1，特征 \(x\) 有年龄、使用时长、投诉次数。最小二乘跑完，你得到一条直线，截距 0.8，系数也都有合理的方向——投诉多的预测值低，使用时间长的预测值高。看起来不错。

然后你看了一眼预测值：好几个样本的拟合数小于 0，还有几个大于 1.5。你试图解释——小于 0 的``更不可能比 0 还要不可能''是什么概念？大于 1 算超过 100\% 的概率？你给输出套上一个 clamp，把超出 \([0,1]\) 的值硬压到端点，但另一个问题立刻浮现：残差散布在 0-1 的两端一定不对称，线性回归的高斯等方差假设根本不成立。

\textbf{你面对的困境是}：一个直线模型试图强行穿过一道概率门。你需要重新回答两件事——数据是怎样生成的，以及什么样的损失函数与这个生成过程保持一致。截断输出只是把问题挡在外面，没有解决它。

\subsection{Logistic 回归：把直线放在一个可以取任意实数的量上}

回到概率的定义。令 \(Y\in\{0,1\}\)，并设

\[
p(x)=P(Y=1\mid x).
\]

概率必须位于 \((0,1)\)，线性组合 \(x^{\trans}\beta\) 却可以取任意实数。我们需要一座桥，把实数值映射到概率区间内。先看 odds——正类概率与负类概率的比值：

\[
\operatorname{odds}(x)=\frac{p(x)}{1-p(x)}.
\]

当 \(p\) 从 0 走到 1，odds 从 0 走到 \(+\infty\)。区间还不够宽：odds 仍被限制在正数一侧。但取对数后，\(\log(\operatorname{odds})\) 可以取遍整条实轴。

这就是模型的切口：把线性结构放在 log-odds 上，而不是概率上：

\[
\log\frac{p(x)}{1-p(x)}=x^{\trans}\beta.
\]

解回概率得到 logistic 函数（sigmoid 的一种参数化）：

\[
p(x)=\sigma(x^{\trans}\beta)
=\frac{1}{1+\exp(-x^{\trans}\beta)}.
\]

这条 S 形曲线在 0 和 1 之间平滑过渡，正是我们需要的变换。现在来解释系数：\(\beta_j\) 不是概率的变化量，而是\textbf{log-odds 的变化量}。若投诉次数的系数为 -0.5，意味着投诉每增加一次，续费 odds 乘以 \(\exp(-0.5)\approx 0.61\)，约降低四成。而且这个倍数对基线概率不同的客户意义不同——一个本来几乎必然续费的客户（\(p=0.95\)，odds=19）和摇摆的客户（\(p=0.5\)，odds=1），同样的 odds 降低转化为概率降幅完全不一样。

\subsection{损失函数：从概率生成到交叉熵}

既然已经规定 \(Y\mid x\sim\operatorname{Bernoulli}(p(x))\)，\(n\) 个独立样本的似然是

\[
\prod_{i=1}^n p_i^{y_i}(1-p_i)^{1-y_i}.
\]

取负对数，这就是我们要最小化的量——二元交叉熵：

\[
L(\beta)
=-\sum_{i=1}^n
\left[
y_i\log p_i+(1-y_i)\log(1-p_i)
\right].
\]

这个损失有两个好处。第一，它是 Bernoulli 负对数似然，不是随便拼凑的。第二，它的 Hessian 具有简洁结构。令 \(p=(p_1,\ldots,p_n)^{\trans}\)，有

\[
\nabla L(\beta)=X^{\trans}(p-y),\qquad
\nabla^2L(\beta)=X^{\trans} WX,
\]

其中 \(W_{ii}=p_i(1-p_i)\ge0\)。Hessian 半正定——线性 Logistic 回归的负对数似然是凸函数。这意味着梯度下降不会把你困在某个局部极小点；任何极小点都是全局的。

但凸性不等于``解一定存在''。下面这个冲突直接击穿了这种安全感。

\subsection{完全分离：当分类完美却无法估计}

假设你手头的数据集里，正负样本可以被一条直线完全分开——训练准确率 100\%，人类直觉中这应该是最简单的情况。但在 Logistic 回归里，这反而是最糟的。

沿分离方向不断放大 \(\beta\)：正类概率趋近 1，负类概率趋近 0，似然持续提高，却没有上限。参数范数 \(\lVert\beta\rVert\) 趋向无穷。有限的最大似然估计不存在——不是算法收敛到错误值，而是根本不存在有限的最优点。

这是个值得反复回味的冲突：分类边界已经完美，概率模型反而无法确定置信程度。分类任务的成功恰好是概率估计的失败。

加入上一节的 Ridge 类惩罚能得到有限解——惩罚项给参数施加了向原点的收缩力，与似然梯度形成平衡。贝叶斯先验也能约束参数的范围。这件事的深层教训是：正则化有时候不是在防止过拟合，而是在修复一个本来就没有良好定义的估计问题。

数值上也必须小心。直接从 sigmoid 输出算 \(\log(1+\exp z)\) 会在 \(|z|\) 很大时溢出：先算 sigmoid 再取对数，可能得到 \(\log 0\)。稳定形式是

\[
\log(1+\exp z)=\max(z,0)+\log(1+\exp(-|z|)),
\]

两个分支都安全。

\subsection{多分类：一个 sigmoid 不够了}

从前面的讨论中你可能已经隐约猜到：如果类别数 \(K>2\)，给每个类单独套一个 sigmoid 会有麻烦——各类概率的和不一定等于 1，而且也没有利用``只能选一个类''的互斥约束。

我们需要的是\textbf{相对分数}，而不是独立分数。给 \(K\) 个类别各自一个 logit：

\[
z_k=w_k^{\trans} x,
\]

然后用 softmax 函数把 \(K\) 个实数归一化成和为 1 的概率向量：

\[
p_k=\frac{\exp(z_k)}{\sum_{j=1}^K\exp(z_j)}.
\]

现在如果一个类的 logit 比其他类高很多，它的概率就接近 1；如果几个类旗鼓相当，概率均匀摊开。对 one-hot 标签 \(y\)，单样本交叉熵为

\[
\ell=-\sum_{k=1}^K y_k\log p_k.
\]

最关键的性质在梯度上：

\[
\nabla_z\ell=p-y.
\]

这是只有 logits 中与真实标签对应的那一维才会出现 ``-1'' 的修正，其余维度全由 softmax 概率承担——梯度的形式与二分类情形如出一辙，只是从标量变成了向量。展开 \(\log p_k\)：

\[
\log p_k = z_k-\log\sum_{j=1}^K e^{z_j}.
\]

第一项负责奖励正确的类，第二项的分母梯度恰好是 softmax 概率 \(p_j\)。

\subsection{参数冗余与数值稳定}

softmax 只依赖相对分数。给所有 logits 加同一个常数 \(c\)：

\[
\frac{\exp(z_k+c)}{\sum_j\exp(z_j+c)}
=\frac{e^c\exp(z_k)}{e^c\sum_j\exp(z_j)}
=p_k.
\]

概率完全不变。参数空间里多了一个方向，沿这个方向的平移不会改变任何预测。你可以固定一个参考类（令其 logit 恒为零），也可以用正则化来挑出代表解。

计算上也利用这种不变性：令 \(m=\max_k z_k\)，则

\[
\log\sum_k e^{z_k}
=m+\log\sum_k e^{z_k-m}.
\]

当 logits 是 \((1000,1001,999)\) 时，直接计算 \(e^{1001}\) 会溢出。减去最大值后变成 \((-1,0,-2)\)，概率不变，计算安全。这条技巧在实现中永远不该省略。

\subsection{多类 vs. 多标签：两种不同的输出结构}

softmax 内置了互斥假设——每个样本只属于一个类别。当一张图片可以同时包含``室内''、``人物''和``夜景''时，标签之间并不互斥。这时应该为每个标签使用独立的 Bernoulli sigmoid，而不是迫使它们在一个共同的分母上竞争。判断的标准不是类别数量，而是这些类别是否在同一组互斥选项内。

\subsection{概率输出和决策是两件事}

模型告诉你 \(P(Y=1\mid x)=0.3\)，不等于``判定为 0''。若漏诊的代价是误报的九倍，即使概率不到 0.5，也可能值得做进一步检查。阈值应该由行动成本、容量约束和基准率共同决定，sigmoid 的 0.5 只是几何中点，没有决策特权。

评估分类器时也需要区分几层概念：

\begin{itemize}
\tightlist
\item
  \textbf{log loss} 衡量整套概率分布的质量——它不仅惩罚错判，也惩罚概率与标签之间过大的距离；
\item
  \textbf{calibration} 问的是：模型预测为 0.7 的样本集合里，是否真的约有 70\% 为正；
\item
  \textbf{ROC 或 PR 曲线} 描述阈值变化时的排序能力——不管阈值取多少，正类得分是否倾向于高于负类；
\item
  \textbf{准确率} 只评价某一个阈值下的硬决策——在类别严重不平衡时，``全部判为多数类''也能拿到很高的准确率。
\end{itemize}

这四者回答的问题互不替代，不能把一个当成另一个的代理。

如果训练时人为过采样正类以平衡数据，你实际上改变了模型所见的 class prior。排序能力可能不受影响，但输出的概率已经偏离真实环境。若部署时的基准率或条件分布继续漂移，历史验证集上的校准也不能无限期保证。模型评估是一项持续的运维工作，不是一次性实验报告。

\subsubsection{从二分类到多分类：一个完整的构建流程}

以三类风险分流为例，整个建模序列可以这样走通：确认类别确实互斥后，选择 softmax 输出层；用减去最大值的 log-sum-exp 实现稳定计算；解析梯度 \(p-y\) 用自动微分或有限差分核对一致；评估时分别报告交叉熵、混淆矩阵和校准曲线。修改输入数据的预处理或给所有 logits 同时加一个常数，softmax 的归一化会自动吸收平移量，预测不变——这是诊断实现正确性的一个快速检查。

若需求变为多标签任务——同一样本可同时属于多个类别——互斥假设不再成立，只需把共享归一化的 softmax 换成各类独立 sigmoid。模型的结构反映了任务的结构：两者匹配时，损失函数所做的优化才对应你想解决的问题。
